\section{Modelo de regresión y diagnóstico}
\label{sec:regresion}

El modelamiento se desarrolló en dos niveles. En primer lugar, se ajustó
una regresión lineal múltiple para estudiar la estructura del monto del
crédito y cumplir con el análisis de regresión y diagnóstico requerido
por el curso. En segundo lugar, se ajustó una regresión logística para
estimar la probabilidad de que un registro sea clasificado como mal
riesgo.

La regresión lineal tiene un papel auxiliar: permite examinar qué
características contractuales y personales están asociadas con el monto
otorgado. No constituye el modelo principal de clasificación crediticia.
La respuesta de riesgo es binaria y se analiza posteriormente mediante
regresión logística.

\subsection{Regresión lineal múltiple para el monto del crédito}
\label{subsec:ols}

\subsubsection{Especificación del modelo}

El monto original presenta una asimetría positiva pronunciada. Por este
motivo, la variable dependiente se expresa en escala logarítmica:

\[
Y_i
=
\log(\variable{monto}_i).
\]

El modelo ajustado fue:

\[
\log(\variable{monto}_i)
=
\beta_0
+
\beta_1\variable{duracion}_i
+
\beta_2\variable{tasa_cuota}_i
+
\beta_3\variable{edad}_i
+
\varepsilon_i.
\]

La especificación incorpora tres dimensiones relevantes de la operación:

\begin{itemize}
    \item la duración contractual del crédito, expresada en meses;

    \item el nivel ordinal de la cuota respecto del ingreso disponible;

    \item la edad del solicitante, expresada en años.
\end{itemize}

Bajo el modelo clásico se supone:

\[
\operatorname{E}(\varepsilon_i\mid X_i)=0,
\]

\[
\operatorname{Var}(\varepsilon_i\mid X_i)=\sigma^2,
\]

y, para la inferencia exacta en muestras finitas:

\[
\varepsilon_i
\sim
\mathcal{N}(0,\sigma^2).
\]

Cada coeficiente representa la variación esperada en el logaritmo del
monto asociada con un aumento de una unidad en el predictor
correspondiente, manteniendo constantes las demás variables incluidas.

Debido a la transformación logarítmica, la variación porcentual asociada
con un coeficiente \(\beta_j\) se obtiene mediante:

\[
100\left[\exp(\beta_j)-1\right].
\]

Esta interpretación es exacta para una variación de una unidad y no debe
confundirse con un efecto causal.

\subsubsection{Ajuste global}

El modelo se estimó sobre las \(1000\) observaciones mediante mínimos
cuadrados ordinarios. La prueba global evalúa:

\[
H_0:
\beta_1=\beta_2=\beta_3=0,
\]

frente a:

\[
H_1:
\text{al menos uno de los coeficientes es distinto de cero}.
\]

El resultado fue:

\[
F(3,996)=386.4,
\qquad
p=\num{2.04e-166}.
\]

Por tanto, se rechaza la hipótesis nula global. Las variables incluidas
aportan información conjunta sobre el logaritmo del monto.

El coeficiente de determinación fue:

\[
R^2=0.538,
\]

mientras que el coeficiente ajustado fue:

\[
R^2_{\mathrm{ajustado}}=0.537.
\]

El modelo explica aproximadamente el \(53.8\,\%\) de la variabilidad
observada en \(\log(\variable{monto})\). El valor restante corresponde a
características no incluidas, variabilidad individual, posibles
relaciones no lineales y error aleatorio.

\subsubsection{Coeficientes estimados}

La Tabla~\ref{tab:coeficientes-ols} presenta los coeficientes clásicos.

\begin{table}[H]
    \centering
    \caption{Coeficientes del modelo OLS para el logaritmo del monto.}
    \label{tab:coeficientes-ols}

    \begin{tabular}{
        @{}
        l
        r
        r
        r
        r
        c
        @{}
    }
        \toprule
        Variable
        & Coeficiente
        & Error estándar
        & Estadístico \(t\)
        & \(p\)
        & IC del \(95\,\%\) \\
        \midrule

        Intercepto
        & \num{7.4860}
        & \num{0.0742}
        & \num{100.915}
        & \(<0.001\)
        & \([7.340,\;7.632]\) \\

        Duración
        & \num{0.0431}
        & \num{0.0014}
        & \num{30.970}
        & \(<0.001\)
        & \([0.040,\;0.046]\) \\

        Tasa de la cuota
        & \num{-0.2476}
        & \num{0.0150}
        & \num{-16.482}
        & \(<0.001\)
        & \([-0.277,\;-0.218]\) \\

        Edad
        & \num{0.0039}
        & \num{0.0015}
        & \num{2.624}
        & \num{0.0088}
        & \([0.001,\;0.007]\) \\

        \bottomrule
    \end{tabular}
\end{table}

Manteniendo constantes las demás variables:

\begin{itemize}
    \item un mes adicional de duración se asocia con un incremento
    aproximado del \(4.4\,\%\) en el monto esperado;

    \item un aumento de una unidad en la escala de
    \variable{tasa_cuota} se asocia con una reducción aproximada del
    \(21.9\,\%\) en el monto esperado;

    \item un año adicional de edad se asocia con un incremento
    aproximado del \(0.4\,\%\) en el monto esperado.
\end{itemize}

La duración presenta la asociación positiva de mayor magnitud. Esto es
coherente con la estructura financiera de una operación: los montos
mayores suelen distribuirse en plazos más largos para reducir la cuota
periódica.

La asociación negativa de \variable{tasa_cuota} también es
económicamente plausible. Una mayor carga relativa de pago puede limitar
el monto que la entidad está dispuesta a conceder. Sin embargo, la
variable está registrada como una escala ordinal de 1 a 4 y no como un
porcentaje exacto. Por ello, el cambio entre niveles se interpreta como
un contraste dentro de la codificación disponible, no como una variación
de un punto porcentual del ingreso.

La edad presenta una asociación estadísticamente significativa, pero de
magnitud reducida. Su importancia práctica es considerablemente menor
que la de la duración y la carga de la cuota.

\begin{hallazgo}
La duración y la carga de la cuota concentran la mayor parte de la
señal explicativa del modelo del monto. La edad aporta información
adicional, pero su asociación condicional es pequeña.
\end{hallazgo}

\subsection{Diagnóstico de multicolinealidad}

La multicolinealidad se evaluó mediante el factor de inflación de la
varianza:

\[
\operatorname{VIF}_j
=
\frac{1}{1-R_j^2},
\]

donde \(R_j^2\) corresponde al coeficiente de determinación obtenido al
regresar el predictor \(X_j\) sobre los demás predictores.

Los resultados se presentan en la
Tabla~\ref{tab:vif-ols}.

\begin{table}[H]
    \centering
    \caption{Factores de inflación de la varianza del modelo OLS.}
    \label{tab:vif-ols}

    \begin{tabular}{lc}
        \toprule
        Variable & VIF \\
        \midrule
        Duración & \num{1.007} \\
        Tasa de la cuota & \num{1.009} \\
        Edad & \num{1.005} \\
        \bottomrule
    \end{tabular}
\end{table}

Los valores se encuentran muy próximos a uno. No se observa una relación
lineal fuerte entre los tres predictores del modelo y, por tanto, no hay
evidencia de multicolinealidad relevante en esta especificación.

Este resultado no contradice la correlación observada entre duración y
monto, porque el monto es la variable dependiente. El VIF evalúa la
relación entre los predictores: duración, tasa de la cuota y edad.

\subsection{Diagnóstico gráfico de los residuos}

La Figura~\ref{fig:diagnosticos-ols} presenta los cuatro gráficos
diagnósticos requeridos.

\begin{figure}[H]
    \centering

    \includegraphics[
        width=\textwidth
    ]{
        figura_06_diagnosticos_modelo_ols.png
    }

    \caption{
        Diagnósticos del modelo OLS para
        \(\log(\variable{monto})\): residuos frente a valores ajustados,
        gráfico Q--Q, gráfico Scale--Location y residuos estandarizados
        frente a leverage.
    }

    \label{fig:diagnosticos-ols}
\end{figure}

El gráfico de residuos frente a valores ajustados no muestra una
curvatura dominante, aunque se observan pequeñas variaciones locales en
la línea suavizada. La dispersión alrededor de cero es razonablemente
estable.

El gráfico Q--Q presenta una alineación adecuada en la región central,
con desviaciones moderadas en las colas. Esto sugiere una aproximación
normal razonable para la mayor parte de los residuos, aunque no una
normalidad estricta.

El gráfico Scale--Location tampoco presenta una forma de embudo marcada.
La línea suavizada muestra variaciones leves, pero no una expansión
sistemática de la varianza a lo largo de los valores ajustados.

Finalmente, el gráfico de residuos frente a leverage permite identificar
observaciones que combinan valores poco habituales de los predictores con
residuos elevados. La distancia de Cook máxima fue:

\[
D_{\max}=\num{0.0476}.
\]

Utilizando el umbral exploratorio:

\[
\frac{4}{n}
=
\frac{4}{1000}
=
0.004,
\]

se identificaron \(52\) observaciones por encima de dicho valor. Este
umbral es deliberadamente sensible y funciona como una señal para
revisión, no como una regla automática de exclusión.

La distancia máxima permanece muy por debajo de uno, por lo que no se
observa una única operación con influencia extrema sobre el ajuste.
Las observaciones señaladas se conservaron porque corresponden a
registros válidos y su eliminación no estaba justificada por un error
documentado.

\subsection{Normalidad de los residuos}

La normalidad se evaluó formalmente mediante la prueba de Shapiro--Wilk:

\[
H_0:
\text{los residuos siguen una distribución normal},
\]

\[
H_1:
\text{los residuos no siguen una distribución normal}.
\]

El resultado fue:

\[
W=\num{0.9951},
\qquad
p=\num{0.00275}.
\]

Al nivel de significancia de \(0.05\), se rechaza la normalidad estricta.
Sin embargo, el valor de \(W\) es cercano a uno y el gráfico Q--Q muestra
que las principales desviaciones se concentran en las colas.

Con \(n=1000\), las pruebas de normalidad tienen capacidad para detectar
desviaciones pequeñas que no necesariamente producen un deterioro
práctico importante de la inferencia. Por esta razón, el resultado se
interpreta junto con los gráficos, el tamaño de la muestra y la
comparación posterior de los errores estándar clásicos y robustos.

\subsection{Homocedasticidad: prueba de Breusch--Pagan}

La constancia de la varianza se evaluó mediante:

\[
H_0:
\operatorname{Var}(\varepsilon_i\mid X_i)=\sigma^2,
\]

\[
H_1:
\operatorname{Var}(\varepsilon_i\mid X_i)
\text{ no es constante}.
\]

La prueba de Breusch--Pagan produjo:

\[
LM=\num{5.888},
\qquad
p=\num{0.11716},
\]

y la versión basada en \(F\):

\[
F=\num{1.967},
\qquad
p=\num{0.11733}.
\]

Como ambos valores \(p\) son superiores a \(0.05\), no se rechaza la
hipótesis de varianza constante. La homocedasticidad constituye una
aproximación razonable para esta especificación.

No rechazar \(H_0\) no demuestra que la varianza sea exactamente
constante. Indica que el test no encontró evidencia suficiente de
heterocedasticidad con los datos y la forma funcional utilizadas.

\subsection{Inferencia robusta HC3}

Como comprobación de sensibilidad, se recalcularon los errores estándar
mediante el estimador robusto HC3. Esta corrección mantiene los
coeficientes y predicciones, pero modifica la matriz de covarianza
utilizada en la inferencia.

\begin{table}[H]
    \centering
    \caption{Comparación entre la inferencia clásica y la corrección
    robusta HC3.}
    \label{tab:hc3-ols}

    \begin{tabular}{
        @{}
        l
        r
        r
        r
        r
        @{}
    }
        \toprule
        Variable
        & Coeficiente
        & EE clásico
        & EE HC3
        & \(p\) HC3 \\
        \midrule

        Intercepto
        & \num{7.4860}
        & \num{0.0742}
        & \num{0.0748}
        & \(<0.001\) \\

        Duración
        & \num{0.0431}
        & \num{0.0014}
        & \num{0.0015}
        & \(<0.001\) \\

        Tasa de la cuota
        & \num{-0.2476}
        & \num{0.0150}
        & \num{0.0156}
        & \(<0.001\) \\

        Edad
        & \num{0.0039}
        & \num{0.0015}
        & \num{0.0016}
        & \num{0.0144} \\

        \bottomrule
    \end{tabular}
\end{table}

Los errores HC3 son ligeramente mayores, pero las conclusiones no
cambian. Duración, tasa de la cuota y edad mantienen asociaciones
estadísticamente significativas.

La estabilidad frente a HC3 respalda que los resultados no dependen
exclusivamente de la matriz de covarianza clásica. La edad continúa
siendo significativa, aunque con menor evidencia que las otras dos
variables.

\subsection{Validación cruzada de cinco particiones}

La capacidad de generalización se evaluó mediante validación cruzada
\(k\)-fold con:

\[
k=5.
\]

La muestra se dividió aleatoriamente en cinco grupos. En cada iteración,
cuatro particiones se utilizaron para ajustar el modelo y la restante
para evaluarlo. La semilla utilizada fue \(\texttt{SEED}=2026\).

El desempeño se midió mediante \(R^2\) y RMSE en la escala de
\(\log(\variable{monto})\).

\begin{table}[H]
    \centering
    \caption{Resultados de validación cruzada del modelo OLS.}
    \label{tab:cv-ols}

    \begin{tabular}{ccc}
        \toprule
        Partición & \(R^2\) & RMSE \\
        \midrule
        1 & \num{0.5321} & \num{0.5776} \\
        2 & \num{0.4972} & \num{0.5292} \\
        3 & \num{0.5943} & \num{0.4959} \\
        4 & \num{0.5439} & \num{0.4861} \\
        5 & \num{0.4901} & \num{0.5561} \\
        \midrule
        Media & \num{0.5315} & \num{0.5290} \\
        Desviación estándar de \(R^2\)
        & \num{0.0418}
        & --- \\
        \bottomrule
    \end{tabular}
\end{table}

El \(R^2\) medio fuera de muestra fue:

\[
\overline{R^2}_{\mathrm{CV}}
=
\num{0.5315}
\pm
\num{0.0418}.
\]

El RMSE medio fue:

\[
\operatorname{RMSE}_{\mathrm{CV}}
=
\num{0.5290}
\]

en la escala logarítmica.

El desempeño de validación es muy próximo al obtenido con la muestra
completa:

\[
R^2_{\mathrm{completo}}=\num{0.5379}.
\]

La diferencia entre ambos valores es reducida, lo que indica que el
modelo mantiene una capacidad explicativa estable en observaciones no
utilizadas durante cada ajuste. No se observa evidencia importante de
sobreajuste dentro del esquema de validación aplicado.

\subsection{Conclusión del modelo lineal}

El modelo OLS proporciona una representación parsimoniosa e
interpretable de la estructura del monto. La duración, la carga de la
cuota y la edad explican conjuntamente cerca del \(54\,\%\) de la
variabilidad de su logaritmo.

La duración presenta la asociación positiva dominante. La tasa de la
cuota mantiene una asociación negativa importante y la edad aporta una
señal positiva pequeña. Los VIF cercanos a uno descartan problemas
relevantes de multicolinealidad entre los predictores seleccionados.

Los diagnósticos no muestran evidencia de heterocedasticidad y la
inferencia robusta HC3 conserva las conclusiones. La prueba de
Shapiro--Wilk detecta una desviación respecto de la normalidad estricta,
pero el comportamiento gráfico, el tamaño de la muestra y la estabilidad
de los errores robustos indican que dicha desviación no modifica de forma
sustantiva la lectura del modelo.

Finalmente, la validación cruzada muestra un desempeño cercano al ajuste
completo. El modelo es útil para comprender la relación entre el plazo,
la capacidad de pago aproximada, la edad y el monto otorgado.

No obstante, este resultado no responde todavía qué solicitantes tienen
mayor probabilidad de mal crédito. Esa pregunta requiere un modelo con
respuesta binaria, desarrollado en el bloque siguiente mediante
regresión logística.

\subsection{Regresión logística para la probabilidad de mal crédito}
\label{subsec:logit}

La segunda parte del modelamiento utiliza como respuesta la variable:

\[
Y_i=\variable{mal_credito}_i,
\]

donde:

\[
Y_i=
\begin{cases}
1, & \text{si el registro corresponde a mal riesgo},\\
0, & \text{si corresponde a buen riesgo}.
\end{cases}
\]

Sea:

\[
p_i
=
\Pr(Y_i=1\mid\mathbf{x}_i)
\]

la probabilidad condicional de que la observación \(i\) sea clasificada
como mal crédito. La regresión logística relaciona esta probabilidad con
los predictores mediante:

\[
\operatorname{logit}(p_i)
=
\log\left(
    \frac{p_i}{1-p_i}
\right)
=
\beta_0+\mathbf{x}_i^{\mathsf{T}}\boldsymbol{\beta}.
\]

De forma equivalente:

\[
p_i
=
\frac{
    \exp\left(
        \beta_0+\mathbf{x}_i^{\mathsf{T}}\boldsymbol{\beta}
    \right)
}{
    1+
    \exp\left(
        \beta_0+\mathbf{x}_i^{\mathsf{T}}\boldsymbol{\beta}
    \right)
}.
\]

El modelo estima probabilidades individuales de mal riesgo. La
transformación de esas probabilidades en una decisión binaria requiere
un umbral, cuya selección se analiza posteriormente.

\subsubsection{Diseño de entrenamiento y prueba}

El conjunto de datos se dividió de forma estratificada en:

\[
70\,\% \text{ para entrenamiento}
\qquad\text{y}\qquad
30\,\% \text{ para prueba}.
\]

La partición resultante fue:

\[
n_{\mathrm{train}}=700,
\qquad
n_{\mathrm{test}}=300.
\]

La tasa de mal crédito se mantuvo en \(30.0\,\%\) en ambos subconjuntos.
Por tanto, el conjunto de entrenamiento contiene \(210\) malos créditos
y el conjunto de prueba contiene \(90\).

La estratificación permite comparar el desempeño sin que una variación
accidental de la prevalencia modifique de manera importante las métricas.
El modelo se ajustó exclusivamente con el conjunto de entrenamiento. El
conjunto de prueba permaneció separado hasta la evaluación final.

La semilla utilizada fue:

\[
\texttt{SEED}=2026.
\]

\subsubsection{Especificación del modelo}

La especificación incorporó variables contractuales, antecedentes
crediticios y señales de estabilidad financiera:

\[
\begin{aligned}
\operatorname{logit}(p_i)
={}&
\beta_0
+\beta_1\variable{duracion}_i
+\beta_2\variable{log_monto}_i
+\beta_3\variable{tasa_cuota}_i
+\beta_4\variable{edad}_i\\
&+\text{estado de cuenta}
+\text{historial crediticio}
+\text{ahorros}\\
&+\text{antigüedad laboral}
+\text{vivienda}
+\text{otros planes de pago}.
\end{aligned}
\]

Las variables categóricas se introdujeron mediante indicadores, dejando
una categoría de referencia en cada grupo. Las referencias fueron:

\begin{itemize}
    \item estado de cuenta: \texttt{A11}, saldo menor que cero;

    \item historial crediticio: \texttt{A30}, sin créditos previos o
    todos pagados;

    \item ahorros: \texttt{A61}, menos de \(100\) DM;

    \item empleo actual: \texttt{A71}, desempleado;

    \item vivienda: \texttt{A151}, alquilada;

    \item otros planes de pago: \texttt{A141}, banco.
\end{itemize}

Los coeficientes se estimaron mediante máxima verosimilitud con
\texttt{statsmodels}. El algoritmo convergió correctamente y el modelo
contiene \(24\) coeficientes, incluido el intercepto.

El pseudo-\(R^2\) de McFadden en entrenamiento fue:

\[
R^2_{\mathrm{McFadden}}
=
0.2202.
\]

Este valor no debe interpretarse como el porcentaje de variabilidad
explicada y no es comparable directamente con el \(R^2\) de una
regresión OLS. Se utiliza como una medida relativa de mejora frente a un
modelo que solo contiene intercepto.

\subsection{Interpretación mediante odds ratios}

En una regresión logística, el coeficiente \(\beta_j\) representa un
cambio en el logaritmo de las \textit{odds}. Para facilitar la
interpretación se utiliza:

\[
\OR_j
=
\exp(\beta_j).
\]

Cuando \(\OR_j>1\), las \textit{odds} de mal crédito aumentan. Cuando
\(\OR_j<1\), disminuyen. En ambos casos, la comparación mantiene
constantes los demás predictores del modelo.

Es importante distinguir \textit{odds} de probabilidad. Un aumento del
\(20\,\%\) en las \textit{odds} no equivale necesariamente a un aumento
de veinte puntos porcentuales en la probabilidad, porque el efecto sobre
la probabilidad depende de su nivel inicial.

\subsubsection{Predictores numéricos}

La Tabla~\ref{tab:or-numericos} presenta los resultados de los
predictores numéricos.

\begin{table}[H]
    \centering
    \caption{Odds ratios de los predictores numéricos del modelo
    logístico.}
    \label{tab:or-numericos}

    \begin{tabular}{
        @{}
        l
        r
        c
        r
        @{}
    }
        \toprule
        Predictor
        & OR
        & IC del \(95\,\%\)
        & \(p\) \\
        \midrule

        Duración
        & \num{1.0390}
        & \([1.0160,\;1.0625]\)
        & \num{0.0008} \\

        Logaritmo del monto
        & \num{1.0658}
        & \([0.7483,\;1.5181]\)
        & \num{0.7238} \\

        Tasa de la cuota
        & \num{1.2283}
        & \([1.0085,\;1.4959]\)
        & \num{0.0409} \\

        Edad
        & \num{0.9972}
        & \([0.9775,\;1.0172]\)
        & \num{0.7798} \\

        \bottomrule
    \end{tabular}
\end{table}

Manteniendo constantes las demás variables, cada mes adicional de
duración se asocia con un aumento aproximado del \(3.9\,\%\) en las
\textit{odds} de mal crédito:

\[
100(1.0390-1)=3.9\,\%.
\]

El intervalo de confianza se mantiene por encima de uno y el valor
\(p\) es inferior a \(0.05\). La duración conserva, por tanto, una
asociación positiva después de controlar por las demás características.

Una unidad adicional en la escala de \variable{tasa_cuota} se asocia con
un incremento aproximado del \(22.8\,\%\) en las \textit{odds}. Esta
interpretación se refiere a un cambio entre niveles de una escala
ordinal, no a un incremento de un punto porcentual del ingreso.

En contraste, \variable{log_monto} deja de ser estadísticamente
significativo:

\[
\OR=1.0658,
\qquad
p=0.7238.
\]

El monto mostró una diferencia marginal entre buenos y malos créditos,
pero no aporta evidencia de una asociación independiente una vez que el
modelo incorpora duración, liquidez, historial y otras variables. Parte
de su señal inicial parece estar compartida con la duración del préstamo.

La edad tampoco presenta evidencia de una asociación condicional dentro
de esta especificación.

\subsubsection{Categorías con evidencia estadística}

La Tabla~\ref{tab:or-categoricos} resume las categorías cuyo intervalo
del \(95\,\%\) no contiene el valor uno.

\begin{table}[H]
    \centering
    \caption{Categorías con asociación estadísticamente significativa en
    el modelo logístico.}
    \label{tab:or-categoricos}

    \footnotesize
    \begin{tabularx}{\textwidth}{
        @{}
        >{\raggedright\arraybackslash}p{3.4cm}
        >{\raggedright\arraybackslash}p{4.1cm}
        >{\centering\arraybackslash}p{1.3cm}
        >{\centering\arraybackslash}p{2.5cm}
        >{\centering\arraybackslash}p{1.4cm}
        @{}
    }
        \toprule
        Variable
        & Categoría frente a referencia
        & OR
        & IC del \(95\,\%\)
        & \(p\) \\
        \midrule

        Ahorros
        & \texttt{A64}: \(1000\) DM o más
          frente a \texttt{A61}
        & \num{0.1617}
        & \([0.0349,\;0.7495]\)
        & \num{0.0199} \\

        Estado de cuenta
        & \texttt{A14}: sin cuenta corriente
          frente a \texttt{A11}
        & \num{0.1803}
        & \([0.1079,\;0.3014]\)
        & \(<0.001\) \\

        Historial crediticio
        & \texttt{A34}: cuenta crítica u otros créditos
          frente a \texttt{A30}
        & \num{0.2413}
        & \([0.0910,\;0.6403]\)
        & \num{0.0043} \\

        Estado de cuenta
        & \texttt{A13}: saldo de \(200\) DM o más
          frente a \texttt{A11}
        & \num{0.3121}
        & \([0.1384,\;0.7038]\)
        & \num{0.0050} \\

        Empleo actual
        & \texttt{A74}: entre cuatro y siete años
          frente a \texttt{A71}
        & \num{0.4037}
        & \([0.1632,\;0.9987]\)
        & \num{0.0497} \\

        Ahorros
        & \texttt{A65}: desconocido o sin cuenta
          frente a \texttt{A61}
        & \num{0.4303}
        & \([0.2421,\;0.7647]\)
        & \num{0.0041} \\

        Vivienda
        & \texttt{A152}: propia
          frente a \texttt{A151}
        & \num{0.5868}
        & \([0.3552,\;0.9694]\)
        & \num{0.0374} \\

        Estado de cuenta
        & \texttt{A12}: saldo entre \(0\) y \(200\) DM
          frente a \texttt{A11}
        & \num{0.6041}
        & \([0.3776,\;0.9665]\)
        & \num{0.0355} \\

        \bottomrule
    \end{tabularx}
\end{table}

El estado de la cuenta corriente presenta un gradiente coherente respecto
de la categoría de saldo negativo. Los perfiles con saldo positivo y los
registros sin una cuenta corriente codificada muestran menores
\textit{odds} que el grupo de referencia.

El resultado de \texttt{A14} no implica que carecer de cuenta corriente
sea, en términos generales, una condición financiera favorable. La
comparación se realiza específicamente frente a solicitantes con saldo
negativo. Además, la categoría histórica puede condensar perfiles
heterogéneos que no pueden separarse con la documentación disponible.

Los niveles altos de ahorro presentan menores \textit{odds} que el grupo
con menos de \(100\) DM. La categoría \texttt{A65} combina ausencia de
cuenta e información desconocida, por lo que su asociación debe
interpretarse con cautela.

La vivienda propia también se asocia con menores \textit{odds} frente a
la vivienda alquilada. La antigüedad laboral de cuatro a siete años
presenta una asociación protectora, aunque su valor \(p=0.0497\) y el
límite superior del intervalo, próximo a uno, indican evidencia débil y
sensible al redondeo.

El resultado del historial \texttt{A34} debe leerse únicamente como un
contraste condicional respecto de \texttt{A30}. La denominación histórica
de las categorías no permite convertir directamente ese coeficiente en
una regla general de aprobación.

\begin{figure}[H]
    \centering

    \includegraphics[
        width=0.92\textwidth
    ]{
        figura_07_odds_ratios_modelo_logistico.png
    }

    \caption{
        Odds ratios e intervalos de confianza del \(95\,\%\) del modelo
        logístico. La línea vertical representa ausencia de asociación
        en la escala de odds ratios.
    }

    \label{fig:odds-ratios-logit}
\end{figure}

La Figura~\ref{fig:odds-ratios-logit} muestra simultáneamente la magnitud,
dirección e incertidumbre de todos los términos. Los intervalos amplios
de algunas categorías reflejan tamaños de grupo reducidos y menor
precisión.

\begin{hallazgo}
La duración mantiene la asociación numérica más consistente con el mal
crédito. El estado de cuenta, los ahorros, el historial, la vivienda y
la carga de la cuota también aportan información condicional. El monto
pierde significancia al controlar por estas variables.
\end{hallazgo}

\subsection{Capacidad discriminativa}

La calidad de un modelo de riesgo no se resume en sus coeficientes. Es
necesario evaluar si las probabilidades estimadas ordenan adecuadamente
los casos de buen y mal crédito en observaciones no utilizadas durante
el ajuste.

La métrica principal independiente del umbral fue el área bajo la curva
ROC:

\[
\AUC_{\mathrm{test}}
=
0.7740.
\]

Este valor puede interpretarse como la probabilidad de que, al elegir al
azar un mal crédito y un buen crédito del conjunto de prueba, el modelo
asigne una puntuación de riesgo mayor al primero.

El resultado indica una capacidad discriminativa claramente superior al
azar, aunque lejos de una separación perfecta. Persisten buenos y malos
créditos con probabilidades similares.

También se calculó el estadístico de Kolmogorov--Smirnov utilizado en
\textit{credit scoring}:

\[
KS
=
\max_t
\left[
    \operatorname{TPR}(t)
    -
    \operatorname{FPR}(t)
\right].
\]

El resultado fue:

\[
KS=0.4810,
\]

alcanzado alrededor del umbral:

\[
t=0.369.
\]

Este resultado muestra una separación apreciable entre las distribuciones
de puntajes de ambas clases. Sin embargo, el umbral que maximiza KS no
tiene por qué minimizar el costo económico de clasificación.

La AUC evalúa capacidad de ordenamiento. No informa por sí sola si las
probabilidades están calibradas, cuál es la rentabilidad de la política
crediticia ni qué umbral debe utilizarse.

\subsection{Validación cruzada del modelo logístico}

La estabilidad se examinó mediante validación cruzada estratificada de
cinco particiones. Dentro de cada partición, las variables categóricas se
codificaron mediante indicadores y las numéricas se estandarizaron usando
exclusivamente la información del conjunto de entrenamiento
correspondiente.

Esta construcción evita que el preprocesamiento utilice información de
la partición de validación.

Los resultados fueron:

\begin{table}[H]
    \centering
    \caption{ROC--AUC por partición en la validación cruzada del modelo
    logístico.}
    \label{tab:cv-logit}

    \begin{tabular}{cc}
        \toprule
        Partición & ROC--AUC \\
        \midrule
        1 & \num{0.7785} \\
        2 & \num{0.7744} \\
        3 & \num{0.8013} \\
        4 & \num{0.7485} \\
        5 & \num{0.7399} \\
        \midrule
        Media & \num{0.7685} \\
        Desviación estándar & \num{0.0221} \\
        \bottomrule
    \end{tabular}
\end{table}

La media de validación fue:

\[
\overline{\AUC}_{\mathrm{CV}}
=
0.7685
\pm
0.0221.
\]

La proximidad entre la AUC de prueba y la media de validación cruzada:

\[
0.7740
\approx
0.7685
\]

indica un desempeño estable entre particiones. No se observa una brecha
importante que sugiera sobreajuste bajo el diseño de evaluación
utilizado.

\subsection{Clasificación con el umbral estándar}

El umbral convencional clasifica una observación como mal riesgo cuando:

\[
\widehat{p}_i
\geq
0.50.
\]

Con este umbral, la matriz de confusión del conjunto de prueba fue:

\[
\begin{array}{c|cc}
 & \text{Predicho bueno} & \text{Predicho malo} \\
\hline
\text{Real bueno} & 180 & 30 \\
\text{Real malo}  & 48  & 42
\end{array}
\]

Por tanto:

\[
TN=180,
\quad
FP=30,
\quad
FN=48,
\quad
TP=42.
\]

Las métricas fueron:

\[
\operatorname{accuracy}
=
0.7400,
\]

\[
\operatorname{precision}
=
0.5833,
\]

\[
\operatorname{recall}
=
0.4667,
\]

\[
F_1
=
0.5185,
\]

\[
\operatorname{especificidad}
=
\frac{180}{180+30}
=
0.8571.
\]

Aunque la exactitud alcanza el \(74.0\,\%\), el modelo identifica menos
de la mitad de los malos créditos. De los \(90\) casos de mal riesgo,
\(48\) serían clasificados como buenos.

Este resultado ilustra por qué la exactitud no puede utilizarse como
único criterio: la política favorece la especificidad, pero conserva
demasiados falsos negativos para una matriz de costos que penaliza ese
error cinco veces más.

\subsection{Selección del umbral bajo costos asimétricos}

La matriz de costos asigna:

\[
C_{\mathrm{FN}}=5,
\qquad
C_{\mathrm{FP}}=1.
\]

Para una observación con probabilidad estimada \(p\), el costo esperado
de clasificarla como buen riesgo es:

\[
C(\text{predecir bueno})
=
C_{\mathrm{FN}}p.
\]

El costo esperado de clasificarla como mal riesgo es:

\[
C(\text{predecir malo})
=
C_{\mathrm{FP}}(1-p).
\]

La decisión de menor costo clasifica como mal riesgo cuando:

\[
C_{\mathrm{FP}}(1-p)
\leq
C_{\mathrm{FN}}p.
\]

Al despejar \(p\):

\[
p
\geq
\frac{
    C_{\mathrm{FP}}
}{
    C_{\mathrm{FP}}+C_{\mathrm{FN}}
}.
\]

Por tanto, el umbral teórico es:

\[
t^{*}
=
\frac{1}{1+5}
=
\frac{1}{6}
\approx
0.1667.
\]

Esta regla supone que los costos representan adecuadamente el problema y
que las probabilidades producidas por el modelo están razonablemente
calibradas. La AUC por sí sola no permite comprobar esta segunda
condición.

El umbral se deriva de la matriz de costos y no de la búsqueda del mejor
resultado en el conjunto de prueba. La curva empírica de costos se
utiliza únicamente como diagnóstico posterior.

\begin{figure}[H]
    \centering

    \includegraphics[
        width=0.88\textwidth
    ]{
        figura_08_costo_vs_umbral.png
    }

    \caption{
        Costo total de clasificación en el conjunto de prueba según el
        umbral. Se muestran el umbral estándar, el umbral teórico
        derivado de la matriz \(5:1\) y el mínimo empírico.
    }

    \label{fig:costo-umbral}
\end{figure}

El costo empírico mínimo se observa alrededor de:

\[
t_{\mathrm{emp}}=0.16,
\]

con un costo de \(157\) unidades. Este valor es prácticamente igual al
umbral teórico \(0.1667\), cuyo costo es \(158\). La diferencia de una
unidad no justifica reemplazar la regla derivada de costos por un umbral
optimizado sobre la misma muestra de prueba.

Seleccionar \(t_{\mathrm{emp}}\) como política final usando las etiquetas
del test introduciría optimismo. En un desarrollo real, cualquier ajuste
empírico del umbral debería efectuarse en validación y confirmarse sobre
un conjunto final independiente.

\subsection{Comparación de las políticas de decisión}

La Tabla~\ref{tab:comparacion-umbrales} compara ambos puntos de
operación.

\begin{table}[H]
    \centering
    \caption{Desempeño del modelo logístico bajo los umbrales
    \(0.50\) y \(1/6\).}
    \label{tab:comparacion-umbrales}

    \footnotesize
    \resizebox{\textwidth}{!}{
        \begin{tabular}{
            @{}
            c
            r
            r
            r
            r
            r
            r
            r
            r
            r
            r
            @{}
        }
            \toprule
            Umbral
            & TP
            & FP
            & FN
            & TN
            & Accuracy
            & Precisión
            & Recall
            & Especificidad
            & \(F_1\)
            & Costo \\
            \midrule

            \(0.50\)
            & 42
            & 30
            & 48
            & 180
            & \num{0.7400}
            & \num{0.5833}
            & \num{0.4667}
            & \num{0.8571}
            & \num{0.5185}
            & 270 \\

            \(1/6\)
            & 80
            & 108
            & 10
            & 102
            & \num{0.6067}
            & \num{0.4255}
            & \num{0.8889}
            & \num{0.4857}
            & \num{0.5755}
            & 158 \\

            \bottomrule
        \end{tabular}
    }
\end{table}

Al reducir el umbral:

\begin{itemize}
    \item los verdaderos positivos aumentan de \(42\) a \(80\);

    \item los falsos negativos disminuyen de \(48\) a \(10\);

    \item el recall aumenta de \(46.7\,\%\) a \(88.9\,\%\);

    \item los falsos positivos aumentan de \(30\) a \(108\);

    \item la precisión disminuye de \(58.3\,\%\) a \(42.6\,\%\);

    \item la especificidad disminuye de \(85.7\,\%\) a \(48.6\,\%\);

    \item la exactitud disminuye de \(74.0\,\%\) a \(60.7\,\%\);

    \item el \(F_1\) aumenta de \(0.5185\) a \(0.5755\).
\end{itemize}

La política basada en costos es más conservadora. Identifica una mayor
proporción de los créditos riesgosos a cambio de marcar también como
riesgosos a más clientes que habrían sido buenos pagadores.

La reducción de falsos negativos es:

\[
48-10=38.
\]

El costo evitado por estos casos es:

\[
38\times5=190.
\]

El incremento de falsos positivos es:

\[
108-30=78,
\]

con un costo adicional de \(78\). El ahorro neto es:

\[
190-78=112.
\]

En consecuencia, el costo total disminuye de:

\[
270
\quad\text{a}\quad
158,
\]

lo que equivale a una reducción de:

\[
\frac{270-158}{270}\times100
=
41.5\,\%.
\]

El costo medio por observación disminuye de \(0.9000\) a \(0.5267\).

\begin{figure}[H]
    \centering

    \includegraphics[
        width=\textwidth
    ]{
        figura_09_roc_y_matrices_confusion.png
    }

    \caption{
        Curva ROC y matrices de confusión del modelo logístico con el
        umbral estándar y el umbral derivado de la matriz de costos.
    }

    \label{fig:roc-matrices}
\end{figure}

La Figura~\ref{fig:roc-matrices} muestra que la curva ROC no cambia al
modificar el umbral: la capacidad de ordenamiento pertenece al modelo.
Lo que cambia es el punto de operación elegido sobre esa curva y, por
tanto, el equilibrio entre falsos positivos y falsos negativos.

\begin{hallazgo}
El umbral \(0.50\) ofrece mayor exactitud y especificidad, pero omite más
de la mitad de los malos créditos. Bajo la matriz de costos \(5:1\), el
umbral \(1/6\) reduce el costo total en aproximadamente \(41.5\,\%\) al
priorizar la detección del riesgo.
\end{hallazgo}

\subsection{Lectura operativa del modelo}

Los resultados permiten diferenciar tres propiedades del sistema:

\begin{enumerate}
    \item \textbf{Asociación.} Los odds ratios identifican variables que
    mantienen una relación condicional con la clase de riesgo.

    \item \textbf{Discriminación.} La AUC y el KS evalúan la capacidad
    para ordenar solicitantes desde menor hasta mayor riesgo.

    \item \textbf{Decisión.} El umbral transforma las probabilidades en
    una política de clasificación y determina el costo de los errores.
\end{enumerate}

Un modelo puede discriminar razonablemente y aun así producir una política
inadecuada si se utiliza un umbral incompatible con los costos del
negocio. De forma similar, un umbral conservador puede reducir pérdidas
por falsos negativos, pero también aumentar rechazos innecesarios.

La matriz \(5:1\) procede de la documentación del dataset y permite
desarrollar el ejercicio de decisión. No representa una estimación
económica completa de una institución. En una aplicación real, el costo
de un falso negativo debería incorporar exposición al incumplimiento,
pérdida dado incumplimiento, recuperación y costos de cobranza. El costo
de un falso positivo debería considerar margen financiero perdido,
costos de adquisición y valor futuro del cliente.

La elección final también requeriría:

\begin{itemize}
    \item evaluar la calibración de las probabilidades;

    \item validar temporalmente el modelo;

    \item comprobar la estabilidad de las variables y categorías;

    \item analizar equidad y posibles efectos adversos;

    \item establecer una zona de revisión manual para probabilidades
    cercanas al punto de corte.
\end{itemize}

Por tanto, el modelo debe considerarse una herramienta de apoyo para
priorizar y revisar solicitudes, no un mecanismo autónomo listo para
aprobar o rechazar créditos.

\subsection{Conclusión del modelamiento}

El modelo OLS mostró que la duración y la carga de la cuota explican una
parte importante de la estructura del monto otorgado, con un desempeño
estable en validación cruzada.

La regresión logística alcanzó una AUC de \(0.7740\) en prueba y una AUC
media de \(0.7685\) en validación cruzada. Esta consistencia indica una
capacidad discriminativa estable y sin evidencia importante de
sobreajuste bajo las particiones evaluadas.

Los factores condicionales más relevantes fueron el estado de la cuenta,
la duración, el historial crediticio, el nivel de ahorros, la vivienda y
la carga de la cuota. El monto mostró una asociación marginal con el
riesgo, pero perdió significancia al controlar simultáneamente por los
demás predictores.

Finalmente, la comparación de umbrales mostró que la política de decisión
no debe basarse exclusivamente en la exactitud. Bajo el costo asimétrico
\(5:1\), reducir el umbral a \(1/6\) aumenta sustancialmente la detección
de malos créditos y reduce el costo total, aunque a cambio de más falsos
positivos.


% 10.1 Modelo OLS auxiliar.
% 10.2 VIF y especificación.
% 10.3 Cuatro diagnósticos canónicos.
% 10.4 Test formal de supuestos.
% 10.5 Validación cruzada k-fold.
% 10.6 Regresión logística para mal riesgo.
% 10.7 Odds ratios e intervalos.
% 10.8 ROC--AUC, precision, recall y matrices de confusión.
% 10.9 Umbral 0,50 frente a umbral basado en costo 5:1.
